% Part: incompleteness
% Chapter: incompleteness-provability
% Section: lob-thm

\documentclass[../../../include/open-logic-section]{subfiles}

\begin{document}

\olfileid[fr]{inc}{inp}{lob}

\olsection{Le théorème de Löb}

L'!!{sentence} de Gödel d'une théorie~$\Th{T}$
est un point fixe de $\lnot\OProv[\Th{T}](y)$ :
c'est un !!{sentence}~$!G$ tel que
\[
\Th{T}\Proves\lnot\OProv[\Th{T}](\gn{!G})
\liff !G.
\]
Il n'est pas démontrable si $\Th{T}$ est
cohérente. En effet, si $\Th{T}\Proves !G$,
alors, d'une part, la condition de dérivabilité
P1 donne
$\Th{T}\Proves\OProv[\Th{T}](\gn{!G})$;
d'autre part, l'équivalence ci-dessus et
$\Th{T}\Proves !G$ donnent
$\Th{T}\Proves\lnot\OProv[\Th{T}](\gn{!G})$.
La théorie serait donc incohérente. Demandons-nous
maintenant ce qu'il en est d'un point fixe de
$\OProv[\Th{T}](y)$, c'est-à-dire d'un
!!{sentence}~$!H$ tel que
\[
\Th{T}\Proves\OProv[\Th{T}](\gn{!H})
\liff !H.
\]
Si $!H$ était démontrable, P1 donnerait
$\Th{T}\Proves\OProv[\Th{T}](\gn{!H})$;
mais l'équivalence ci-dessus, appliquée à
$\Th{T}\Proves !H$, donnerait la même conclusion
par modus ponens. Nous n'obtenons donc pas de
contradiction par l'argument utilisé pour
l'!!{sentence} de Gödel. Cela ne prouve évidemment
pas que $\Th{T}$ \emph{démontre}~$!H$.

Pour avancer, généralisons quelque peu la question.
Le sens de gauche à droite de l'équivalence du
point fixe,
$\OProv[\Th{T}](\gn{!H})\lif !H$,
est une instance d'un schéma général appelé
\emph{principe de réflexion} :
$\OProv[\Th{T}](\gn{!A})\lif !A$.
Ce nom évoque la capacité de~$\Th{T}$ à
« réfléchir » sur ce qu'elle peut démontrer.
En substance, le principe dit : « Si $\Th{T}$
peut démontrer~$!A$, alors $!A$ est vrai », pour
tout~$!A$. Toutes ses instances ne sont vraies
que pour une théorie correcte, ce qui laisse penser
qu'une théorie ne les démontre pas toutes en
général. Quelles instances une théorie suffisamment
forte et satisfaisant les conditions de dérivabilité
peut-elle donc démontrer? Assurément celles où
$!A$ est lui-même démontrable. Le résultat suivant
montre qu'il n'y en a pas d'autres.

\begin{thm}\ollabel{thm:lob}
Soit $\Th{T}$ une théorie !!{axiomatizable}
étendant~$\Th{Q}$, et supposons que la
!!{formula}~$\OProv[\Th{T}](y)$ satisfasse les
conditions P1--P3 de la \olref[prc]{sec}. Pour tout
!!{sentence}~$!A$, si
\[
\Th{T}\Proves\OProv[\Th{T}](\gn{!A})\lif !A,
\]
alors $\Th{T}\Proves !A$.
\end{thm}

Autrement dit, si $\Th{T}\Proves/ !A$, alors
$\Th{T}\Proves/
\OProv[\Th{T}](\gn{!A})\lif !A$.
C'est le théorème de Löb.

\begin{explain}
Voici une intuition pour la preuve du théorème de
Löb, sous la forme d'un raisonnement ingénieux
concluant à l'existence du Père Noël. Si cette
conclusion ne vous plaît pas, remplacez-la par
n'importe quelle autre. L'argument est le suivant.
\begin{enumerate}
\item Soit $X$ la phrase : « Si $X$ est vraie,
  alors le Père Noël existe. »
\item Supposons que $X$ soit vraie.
\item Ce qu'elle dit vaut donc : si $X$ est vraie,
  alors le Père Noël existe.
\item Puisque nous supposons $X$ vraie, (2) et~(3)
  donnent, par modus ponens, que le Père Noël existe.
\item Nous avons déduit (4), « le Père Noël existe »,
  de l'hypothèse~(2), « $X$ est vraie ». Par preuve
  conditionnelle, nous avons donc montré :
  « Si $X$ est vraie, alors le Père Noël existe. »
\item Or, il s'agit précisément de la phrase~$X$.
  Nous avons donc montré que $X$ est vraie.
\item Mais alors, par l'argument (2)--(4), le
  Père Noël existe.
\end{enumerate}
Ce raisonnement est une forme du paradoxe de Curry. Son
étape (3) traite la vérité de la phrase comme une autorisation
d'affirmer ce qu'elle dit. Dans l'arithmétique, passer de la
démontrabilité d'un énoncé à cet énoncé exige précisément une
instance du principe de réflexion. Pour rendre l'idée formelle,
on applique le lemme du point fixe à la !!{formula}
$\OProv[\Th{T}](y)\lif !A$. Le point fixe joue le rôle de~$X$.
La preuve utilise P1--P3 et suppose l'instance de réflexion
pour le seul énoncé~$!A$.
\end{explain}

\begin{proof}[Preuve du \olref{thm:lob}]
Supposons que $!A$ soit un !!{sentence} et que
\[
\Th{T}\Proves\OProv[\Th{T}](\gn{!A})\lif !A.
\]
Posons $!B(y)=\OProv[\Th{T}](y)\lif !A$.
Le lemme du point fixe donne un !!{sentence}~$!D$
tel que
\[
\Th{T}\Proves !D\liff !B(\gn{!D}).
\]
Chacune des lignes suivantes est démontrable
dans~$\Th{T}$ :
\begin{align}
  & !D \liff (\OProv[\Th{T}](\gn{!D}) \lif !A) \ollabel{L-1}\\
  & \qquad \text{$!D$ est un point fixe de~$!B(y)$}\notag \\
  & !D \lif (\OProv[\Th{T}](\gn{!D}) \lif !A) \ollabel{L-2}\\
  & \qquad\text{par \olref{L-1}}\notag\\
  & \OProv[\Th{T}](\gn{!D \lif (\OProv[\Th{T}](\gn{!D}) \lif !A)}) \ollabel{L-3}\\
  & \qquad\text{par \olref{L-2} et P1}\notag \\
  & \OProv[\Th{T}](\gn{!D}) \lif \OProv[\Th{T}](\gn{\OProv[\Th{T}](\gn{!D}) \lif !A})
  \ollabel{L-4}\\
  &\qquad \text{par \olref{L-3} et P2}\notag \\
  & \OProv[\Th{T}](\gn{!D}) \lif (\OProv[\Th{T}](\gn{\OProv[\Th{T}](\gn{!D})}) \lif \OProv[\Th{T}](\gn{!A})) \ollabel{L-5}\\
  &\qquad \text{par \olref{L-4} et P2}\notag\\
  & \OProv[\Th{T}](\gn{!D}) \lif \OProv[\Th{T}](\gn{\OProv[\Th{T}](\gn{!D})}) \ollabel{L-6}\\
  & \qquad\text{par P3}\notag\\
  & \OProv[\Th{T}](\gn{!D}) \lif \OProv[\Th{T}](\gn{!A}) \ollabel{L-7}\\
  &\qquad\text{par \olref{L-5} et \olref{L-6}}\notag\\
  & \OProv[\Th{T}](\gn{!A}) \lif !A \ollabel{L-8}\\
  &\qquad\text{par hypothèse}\notag\\
  & \OProv[\Th{T}](\gn{!D}) \lif !A \ollabel{L-9}\\
  &\qquad\text{par \olref{L-7} et \olref{L-8}}\notag\\
  & (\OProv[\Th{T}](\gn{!D}) \lif !A) \lif !D \ollabel{L-10}\\
  & \qquad \text{par \olref{L-1}}\notag \\
  & !D \ollabel{L-11}\\
  & \qquad\text{par \olref{L-9} et \olref{L-10}}\notag \\
  & \OProv[\Th{T}](\gn{!D}) \ollabel{L-12}\\
  & \qquad\text{par \olref{L-11} et P1}\notag \\
  & !A \qquad\qquad\text{par \olref{L-9} et \olref{L-12}}\notag
\end{align}
\end{proof}

Le théorème de Löb donne une courte preuve du
second théorème d'incomplétude pour les théories
dont le prédicat de dérivabilité satisfait P1--P3 :
si $\Th{T}\Proves
\OProv[\Th{T}](\gn{\lfalse})\lif\lfalse$,
alors $\Th{T}\Proves\lfalse$.
Si $\Th{T}$ est cohérente, elle ne démontre pas
$\lfalse$; elle ne démontre donc pas non plus
$\OProv[\Th{T}](\gn{\lfalse})\lif\lfalse$,
c'est-à-dire $\OCon[\Th{T}]$.
Le même théorème montre que le point fixe~$!H$
de $\OProv[\Th{T}](x)$ est démontrable. En effet,
\begin{align*}
  \Th{T} & \Proves \OProv[\Th{T}](\gn{!H}) \liff !H\\
  \intertext{et donc, en particulier,}
    \Th{T} & \Proves \OProv[\Th{T}](\gn{!H}) \lif !H.
\end{align*}
Par le théorème de Löb, $\Th{T}\Proves !H$.

\begin{prob}
Soit $\Th{T}$ une théorie axiomatisée par un
ensemble calculable et soit $\OProv[\Th{T}]$
un prédicat de dérivabilité pour~$\Th{T}$.
Considérez les quatre énoncés suivants :
\begin{enumerate}
\item Si $\Th{T}\Proves !A$, alors
  $\Th{T}\Proves\OProv[\Th{T}](\gn{!A})$.
\item $\Th{T}\Proves !A\lif
  \OProv[\Th{T}](\gn{!A})$.
\item Si $\Th{T}\Proves
  \OProv[\Th{T}](\gn{!A})$, alors
  $\Th{T}\Proves !A$.
\item $\Th{T}\Proves
  \OProv[\Th{T}](\gn{!A})\lif !A$.
\end{enumerate}
Sous quelles conditions chacun de ces énoncés
est-il vrai?
\end{prob}

\end{document}
